\section*{Summary}
\addcontentsline{toc}{section}{Summary}

\textbf{The question.} The census of quantum-classical gaps in this repository counts \emph{classes} of
structures that are proven to have the same quantum model set. Those classes are built from rules, and every rule
has to be valid for quantum sources. The rules were taken over from the classical theory, where a latent is a
random variable with a value. A quantum source has no value, only shares that are measured, and two of the three
classical rules use the value.

\textbf{The answer in one table.} A tick means proven; a cross means refuted by a counterexample in this note.

\begin{center}
\small
\renewcommand{\arraystretch}{1.25}
\begin{tabular}{@{}>{\raggedright\arraybackslash}p{3.0cm}>{\raggedright\arraybackslash}p{5.8cm}cc>{\raggedright\arraybackslash}p{2.5cm}@{}}
\toprule
\textbf{Rule} & \textbf{Hypothesis, in one line} & \textbf{Classical} & \textbf{Quantum} & \textbf{Where}\\
\midrule
Structural dominance & fewer edges or smaller sources allow fewer distributions
  & \checkmark & \checkmark & \cref{lem:monotone}\\
Face splitting, closed families & every source that touches the split-off part $C$ is split; the rest $D$ sees
  the outputs of $C$ and its visible parents & \checkmark & \checkmark & \cref{thm:split}\\
Face splitting, moderate and strong & as above, but a node of $C$ may also sit in an unsplit source that
  contains $D$ & \checkmark & $\boldsymbol{\times}$ & \cref{thm:sixnode}\\
Edge dropping, Rule A & the origin has no visible parent and a single source, which contains the target
  & \checkmark & \checkmark & \cref{thm:ruleA}\\
Edge dropping, Rule B & every other member of the origin's sources is premeasurable; the target sees the
  origin's inputs & \checkmark & \checkmark & \cref{thm:ruleB}\\
Edge dropping, the classical rule & the target sees the origin's visible parents and shares all its sources
  & \checkmark & $\boldsymbol{\times}$ & \cref{thm:hlp-settings,thm:hlp-four}\\
\bottomrule
\end{tabular}
\end{center}

\begin{figure}[H]
\centering
\begin{tikzpicture}[mdag,scale=0.7]
  \tikzset{panel/.style={draw=#1, line width=1pt, rounded corners=3pt, inner sep=6pt},
           ptitle/.style={font=\scriptsize\bfseries, text=#1, anchor=south west, inner sep=1.5pt}}
  % ---------- Rule 1
  \begin{scope}[shift={(0,0)}]
    \vn[Cn]{a}{0,1.3}{$a$} \vn[Cn]{b}{1.5,1.3}{$b$} \vn[Dn]{d}{0.75,0}{$d$} \qs[splitsrc]{q}{0.75,0.85}
    \feed[splitedge]{q}{a,b,d} \de{a}{d} \de{b}{d}
    \node[rel] at (2.4,0.65) {\equivsign};
    \begin{scope}[shift={(3.3,0)}]
      \vn[Cn]{a2}{0,1.3}{$a$} \vn[Cn]{b2}{1.5,1.3}{$b$} \vn[Dn]{d2}{0.75,0}{$d$} \qs{q2}{0.75,2.05}
      \feed{q2}{a2,b2} \de{a2}{d2} \de{b2}{d2}
    \end{scope}
    \begin{scope}[on background layer]\node[panel=oigreen!85!black, fit=(a)(d)(q2)(b2)(d2)] (p1) {};\end{scope}
    \node[ptitle=oigreen!60!black] at (p1.north west) {\checkmark\ Rule 1: split a source};
  \end{scope}
  % ---------- Rule A
  \begin{scope}[shift={(6.9,0)}]
    \vn[pre]{a}{0,0.2}{$a$} \vn{b}{1.6,0.2}{$b$} \qs{q}{0.8,1.2} \feed{q}{a,b} \de[dropedge]{a}{b}
    \node[rel] at (2.5,0.65) {\equivsign};
    \begin{scope}[shift={(3.4,0)}]
      \vn[pre]{a2}{0,0.2}{$a$} \vn{b2}{1.6,0.2}{$b$} \qs{q2}{0.8,1.2} \feed{q2}{a2,b2}
    \end{scope}
    \node (top) at (0.8,2.15) {}; \node (bot) at (0.8,-0.3) {};
    \begin{scope}[on background layer]\node[panel=oigreen!85!black, fit=(a)(b2)(top)(bot)] (p2) {};\end{scope}
    \node[ptitle=oigreen!60!black] at (p2.north west) {\checkmark\ Rule 2A: drop an edge};
  \end{scope}
  % ---------- Rule B
  \begin{scope}[shift={(13.9,0)}]
    \vn{y}{0.8,1.8}{$y$} \vn{a}{0,0.45}{$a$} \vn{b}{1.6,0.45}{$b$} \qs{q}{0.8,-0.25}
    \feed{q}{a,b} \de{y}{a} \de{y}{b} \de[dropedge]{a}{b}
    \node[rel] at (2.5,0.65) {\equivsign};
    \begin{scope}[shift={(3.4,0)}]
      \vn{y2}{0.8,1.8}{$y$} \vn{a2}{0,0.45}{$a$} \vn{b2}{1.6,0.45}{$b$} \qs{q2}{0.8,-0.25}
      \feed{q2}{a2,b2} \de{y2}{a2} \de{y2}{b2}
    \end{scope}
    \node (top) at (0.8,2.15) {};
    \begin{scope}[on background layer]\node[panel=oigreen!85!black, fit=(a)(b2)(top)(q)] (p3) {};\end{scope}
    \node[ptitle=oigreen!60!black] at (p3.north west) {\checkmark\ Rule 2B: drop an edge};
  \end{scope}
  % ---------- classical edge rule fails
  \begin{scope}[shift={(0.2,-4.9)}]
    \vn{s}{0.9,2.3}{$S$} \vn{a}{0,1.0}{$A$} \vn{b}{1.8,1.0}{$B$} \vn{c}{0.9,-0.4}{$C$} \vn{t}{2.4,-0.4}{$T$}
    \qs{q}{0.9,0.4} \feed{q}{a,b,c} \de{s}{a} \de{s}{b} \de{t}{c} \de{a}{b}
    \node[rel] at (3.35,1.0) {\nequivsign};
    \begin{scope}[shift={(4.3,0)}]
      \vn{s2}{0.9,2.3}{$S$} \vn{a2}{0,1.0}{$A$} \vn{b2}{1.8,1.0}{$B$} \vn{c2}{0.9,-0.4}{$C$} \vn{t2}{2.4,-0.4}{$T$}
      \qs{q2}{0.9,0.4} \feed{q2}{a2,b2,c2} \de{s2}{a2} \de{s2}{b2} \de{t2}{c2}
    \end{scope}
    \begin{scope}[on background layer]\node[panel=oired, fit=(s)(a)(c)(t2)(s2)] (p4) {};\end{scope}
    \node[ptitle=oired] at (p4.north west) {$\boldsymbol{\times}$ the classical edge rule};
  \end{scope}
  % ---------- moderate face splitting fails
  \begin{scope}[shift={(10.0,-4.9)}]
    \vn{s}{-0.5,2.4}{$S$} \vn[Cn]{A}{0,1.3}{$A$} \vn[Cn]{R}{1.8,1.3}{$R$} \vn[Dn]{D}{0.9,-0.2}{$D$} \vn{E}{2.9,-0.2}{$E$} \vn{t}{3.6,1.0}{$T$}
    \qs[splitsrc]{f}{0.9,0.85} \qs{g}{1.9,0.25}
    \feed[splitedge]{f}{A,R,D} \feed{g}{R,D,E} \de{A}{D} \de{R}{D} \de{s}{A} \de{t}{E}
    \draw[edge] (s) to[out=-150,in=165,looseness=1.3] (D);
    \node[rel] at (4.55,1.0) {\nequivsign};
    \begin{scope}[shift={(6.7,0)}]
      \vn{s2}{-0.5,2.4}{$S$} \vn[Cn]{A2}{0,1.3}{$A$} \vn[Cn]{R2}{1.8,1.3}{$R$} \vn[Dn]{D2}{0.9,-0.2}{$D$} \vn{E2}{2.9,-0.2}{$E$} \vn{t2}{3.6,1.0}{$T$}
      \qs{f2}{0.9,2.05} \qs{g2}{1.9,0.25}
      \feed{f2}{A2,R2} \feed{g2}{R2,D2,E2} \de{A2}{D2} \de{R2}{D2} \de{s2}{A2} \de{t2}{E2}
      \draw[edge] (s2) to[out=-150,in=165,looseness=1.3] (D2);
    \end{scope}
    \node (left) at (-1.5,1.0) {};
    \begin{scope}[on background layer]\node[panel=oired, fit=(left)(s)(D)(t2)(s2)(D2)] (p5) {};\end{scope}
    \node[ptitle=oired] at (p5.north west) {$\boldsymbol{\times}$ moderate and strong face splitting};
  \end{scope}
\end{tikzpicture}
\caption{The note at a glance. Top: the three rules that hold for quantum sources, each in its smallest instance
(\cref{sec:split,sec:edge}). Bottom: two pairs that the classical rules declare equivalent and that have
different quantum model sets (\cref{sec:notquantum}).}
\label{fig:glance}
\end{figure}

\textbf{Three things to remember.}
\begin{enumerate}[leftmargin=*]
\item \emph{Every quantum proof here is one of two moves.} Either a source measures a share itself, because the
  measurement was fixed anyway (\emph{pre-measurement}), or a source prepares in advance one state for every value
  that its receivers will later see, and they pick the right one (\emph{the catalogue}). Both moves only ever
  condition on \emph{visible} values.
\item \emph{The classical proofs that fail condition on a latent's value.} The classical edge rule lets the
  shared latent carry the origin's noise; the classical moderate and strong face-splitting rules sample the new
  latent of $D$ from a posterior given the other latents that $C$ sees. Neither has a quantum counterpart, and in
  both cases the conclusion itself fails (\cref{sec:notquantum}).
\item \emph{The code now follows the proofs.} The quantum metagraph of the repository uses face splitting only for
  closed families (\code{mDAG.quantum\_splittable\_families}); the classical metagraph keeps the strong rule. On
  four visible nodes this gives 446 proven quantum classes instead of 438. No conclusion about gaps changes: the
  same 46 classes remain open (\cref{sec:buys}).
\end{enumerate}

\textbf{How to read the pictures.} \Cref{fig:legend} shows every graphical convention used below.

\begin{figure}[H]
\centering
\begin{tikzpicture}[mdag]
  % a small structure
  \vn{x}{0,0}{$x$} \vn{a}{1.6,0}{$a$} \vn{b}{3.6,0}{$b$}
  \qs{q}{2.6,0.95}
  \de{x}{a} \de{a}{b} \feed{q}{a,b}
  \node[lab,anchor=north] at (1.8,-0.55) {\str{x→a, a→b; Q\{a,b\}}};
  % legend
  \begin{scope}[shift={(5.6,0.85)}]
    \vn{l1}{0,0}{$v$} \node[lab,anchor=west] at (0.5,0) {visible node: its outcome is classical};
    \draw[edge] (-0.3,-0.65) -- (0.3,-0.65); \node[lab,anchor=west] at (0.5,-0.65) {visible edge: the outcome is sent along it};
    \qs{l2}{0,-1.3} \draw[qedge] (l2) -- ++(0.45,-0.3); \node[lab,anchor=west] at (0.5,-1.3) {quantum source and the shares it sends};
    \draw[cedge] (-0.3,-1.95) -- (0.3,-1.95); \node[lab,anchor=west] at (0.5,-1.95) {classical data handed out by a source};
  \end{scope}
  \begin{scope}[shift={(12.4,0.85)}]
    \vn[Cn]{m1}{0,0}{$c$} \node[lab,anchor=west] at (0.5,0) {part $C$, split off};
    \vn[Dn]{m2}{0,-0.65}{$d$} \node[lab,anchor=west] at (0.5,-0.65) {part $D$, which sees $C$};
    \vn[pre]{m3}{0,-1.3}{$p$} \node[lab,anchor=west] at (0.5,-1.3) {premeasurable node};
    \draw[dropedge] (-0.3,-1.95) -- (0.3,-1.95); \node[lab,anchor=west] at (0.5,-1.95) {edge that is dropped};
  \end{scope}
\end{tikzpicture}
\caption{Conventions. Structures are also written in the repository's notation: the visible edges, then the
sources, \str{Q\{\dots\}} for a quantum source. An orange source is the one a rule acts on.}
\label{fig:legend}
\end{figure}
